\documentclass[12pt]{article}
\usepackage[utf8]{inputenc}
\usepackage[T1]{fontenc}
\usepackage{hyperref}
\title{Barbaram Language Resources for K-12 Curriculum Development}
\author{Compiled by Claude Sonnet 5 (claude-sonnet-5) for Andr\'as Kornai}
\date{}
\begin{document}
\maketitle
\section{Introduction}
Barbaram (usually spelled Mbabaram; ISO 639-3 vmb, Glottocode mbab1239) is an extinct Pama-Nyungan language of north Queensland, Australia. Sources disagree on when the last speaker died (1972 or 1979). This report found no ready-made school materials; only reference pages were located by search.
\section{Primary Educational Resources}
No primers, textbooks or school programmes for Barbaram were found in search results.
\section{Linguistic and Reference Materials}
Overview of the language and its classification: \url{https://en.wikipedia.org/wiki/Mbabaram_language}. Background on the people: \url{https://en.wikipedia.org/wiki/Mbabaram_people}. A Rosetta Panglossia entry: \url{https://rosettapanglossia.longnow.org/index.php?title=Vmb}. Search summaries state that most of what is known derives from R.\,M.\,W. Dixon's fieldwork with the last speaker, and that recordings are held at AIATSIS; no URL for these was seen, so they are not cited.
\section{Religious and Community Resources}
None found.
\section{Digital and Audio Resources}
None found beyond the AIATSIS recordings mentioned above, which could not be linked.
\section{Curriculum Development Context}
The language is extinct, so a heritage curriculum would have to be built from Dixon's documentation and AIATSIS holdings, in consultation with the Mbabaram community. Little is available online.
\end{document}